\documentclass[10pt]{article}
\usepackage[T1]{fontenc}
\usepackage{lmodern}
\usepackage[a4paper,margin=28mm]{geometry}
\usepackage{amsmath,amssymb,amsthm}
\usepackage[hidelinks]{hyperref}
\hypersetup{pdftitle={On a counting conjecture for binary rich words},pdfauthor={Max Glass}}
\newtheorem{theorem}{Theorem}
\newtheorem{lemma}[theorem]{Lemma}
\newtheorem{proposition}[theorem]{Proposition}
\newtheorem{corollary}[theorem]{Corollary}
\newcommand{\Pal}{\mathcal P}
\newcommand{\Alphab}{\operatorname{alph}}
\title{On a counting conjecture for binary rich words}
\author{Max Glass}
\date{November 6, 2026}

\begin{document}
\maketitle
\vspace{-2em}
\begin{abstract}
Let $R_2(n)$ denote the number of binary rich words of length $n$.
We prove that
\[
 \limsup_{n\to\infty}\frac{\log R_2(n)}{\sqrt n\log n}=\infty.
\]
Consequently, $R_2(n)$ is not $O(n^{C\sqrt n})$ for any fixed $C>0$,
disproving the conjectured upper bound
$O((n/g(n))^{\sqrt n})$ with $g(n)\to\infty$.
The proof combines independent permutation families through an ordered
concatenation construction and a binary gap encoding.
\end{abstract}

\section{The result}

Let $R_2(n)$ count rich words of length $n$ over $\{0,1\}$.
Guo, Shallit, and Shur~\cite[Conjecture~1]{GSSpre} conjectured the bound
\begin{equation}\label{eq:conjecture}
 R_2(n)=O\!\left(\left(\frac{n}{g(n)}\right)^{\sqrt n}\right)
 \qquad\text{for some positive }g(n)\longrightarrow\infty.
\end{equation}
This formulation also appears in~\cite[Section~1]{Ruk}.
Using run-length constructions, Guo, Shallit, and Shur~\cite{GSS}
obtained a lower bound exponential in $\sqrt n$.

\begin{theorem}\label{thm:main}
The binary rich-word counting function satisfies
\[
 \limsup_{n\to\infty}\frac{\log R_2(n)}{\sqrt n\log n}=\infty.
\]
\end{theorem}

\begin{corollary}\label{cor:conjecture}
For every $C>0$, $R_2(n)\ne O(n^{C\sqrt n})$.
In particular, \eqref{eq:conjecture} is false.
\end{corollary}
\begin{proof}
An eventual bound $R_2(n)\le K n^{C\sqrt n}$ would give a limsup at most $C$.
Since $g(n)\ge1$ eventually, \eqref{eq:conjecture} would imply such a
bound with $C=1$.
\end{proof}

The construction starts with $r!$ rich blocks of length $O(r^2)$.
To multiply the choices, we place successive blocks on disjoint ordered
pairs of letters and then encode the resulting word by binary gaps.
The encoding charges $O(jr_j^2)$ for block $j$, suggesting parameters
$r_j$ proportional to $1/j$. We implement this choice dyadically:
$2^t$ blocks of parameter $2s\,2^{k-t}$ at level $t$. Each level has
length cost $O(s^2 4^k)$ and contributes at least $s2^k\log s$ to the
logarithm of the count. Summing over $k+1$ levels gives a lower bound of
order $\sqrt{k+1}$ for the normalized logarithm, with $s\to\infty$
first. All logarithms are natural.

\section{Ordered concatenation}

We use the standard facts that a word $w$ has at most $|w|+1$ distinct
palindromic factors and that every factor of a rich word is rich
\cite[Proposition~2.1 and Corollary~2.5]{GJW}.
Write $\Pal(w)$ for the set of palindromic factors of $w$, including
$\varepsilon$, and $D(w)=|w|+1-|\Pal(w)|$ for its palindromic defect.
For a set $S$ of letters, let $|w|_S$ count the positions of $w$ carrying
letters of $S$, and put
\[
 \begin{aligned}
 I_S(w)&=\{p\in\Pal(w):apb\text{ is a factor of }w
                    \text{ for some }a,b\in S\},\\
 \sigma_S(w)&=\max\{|w|_S-1,0\}-|I_S(w)|.
 \end{aligned}
\]
The two endpoint letters need not coincide. We abbreviate $I_{\{b\}}$
to $I_b$. If $w$ is rich and $b$ occurs in $w$, then
\begin{equation}\label{eq:singleton}
 |I_b(w)|=|w|_b-1.
\end{equation}
Each position carrying $b$ first introduces exactly one palindrome ending
in $b$; these are $b$ and the words $bpb$ with $p\in I_b(w)$
\cite[proof of Proposition~2.7]{GJW}.

If all letters are marked and $|w|\ge2$, then
$I_{\Alphab(w)}(w)$ is the set of palindromes in $w$ with its first and last
letters deleted. Thus $\sigma_{\Alphab(w)}(w)=0$ for rich $w$, also when
$|w|\le1$.

For a word $d=a_1\cdots a_m$ of nonnegative integers, define
\[
 G(d)=0\,1^{a_1}0\,1^{a_2}0\cdots1^{a_m}0,
 \qquad G(\varepsilon)=0.
\]
For nonempty $d$, with $M=\max_i a_i$ and $S_h=\{a:a\ge h\}$, one has
\begin{equation}\label{eq:defect}
 D(G(d))=D(d)+\sum_{h=1}^{M}\sigma_{S_h}(d).
\end{equation}
To see this, the nonempty zero-ended palindromes are exactly $G(p)$,
$p\in\Pal(d)$. The one-ended palindromes are $1,\ldots,1^M$ and
$1^hG(p)1^h$ with $p\in I_{S_h}(d)$. Parsing at zeros makes these
classes disjoint and exhaustive. Hence
\[
 |\Pal(G(d))|=1+|\Pal(d)|+M+\sum_{h=1}^{M}|I_{S_h}(d)|.
\]
Now use $|G(d)|=1+m+\sum_i a_i$ and
$\sum_{h=1}^{M}|d|_{S_h}=\sum_i a_i$.

\begin{lemma}\label{lem:concat}
Let $q\ge1$ and let $z_1,\ldots,z_q$ be nonempty binary words ending in $1$, with
$z_j1$ rich for every $j$. Relabel $0,1$ in $z_j$ as $2j-2,2j-1$,
respectively, obtaining $B_j$. Then $G(B_1\cdots B_q)$ is rich.
\end{lemma}
\begin{proof}
Put $d=B_1\cdots B_q$. Each block is rich, and their alphabets are
pairwise disjoint. No nonempty palindrome crosses a block boundary,
so $d$ is rich. We show that every term in \eqref{eq:defect} vanishes.

At an even threshold $h=2j-2$, the marked positions form the entire
rich suffix $B_j\cdots B_q$, so the marked deficit is zero.
At an odd threshold $h=2j-1$, put $U=B_j$, $b=2j-1$, and
$V=B_{j+1}\cdots B_q$. Only $b$ is marked in $U$, and all letters
are marked in $V$. If $V$ is empty, use \eqref{eq:singleton}.
Otherwise, writing $V^-$ for $V$ with its last letter deleted,
\begin{equation}\label{eq:boundary}
 I_{\{b\}\cup\Alphab(V)}(UV)=I_b(Ub)\cup\Pal(V^-).
\end{equation}
For an interior contained in $U$, the first letter of $V$ can replace
the auxiliary final $b$ in $Ub$. For an interior contained in $V$,
the last letter of $U$ supplies a left endpoint. No nonempty palindrome
crosses the disjoint-alphabet boundary. These observations give
\eqref{eq:boundary}. Its two sets intersect only in $\varepsilon$,
which belongs to both since $Ub$ ends in $bb$.
As $Ub$ and $V^-$ are rich, the union has
$|U|_b+|V|-1$ elements. The marked deficit is therefore zero.
Earlier blocks contain no marked positions and do not affect these
interiors. Equation~\eqref{eq:defect} completes the proof.
\end{proof}

\section{A product lower bound}

\begin{lemma}\label{lem:permutation}
For each $r\ge1$ there are $r!$ distinct binary words $z$ ending in $1$
such that $z1$ is rich and
\begin{equation}\label{eq:composition}
 |z|=L_r:=\frac{5r^2+9r+4}{2},\qquad |z|_1=(r+1)^2.
\end{equation}
\end{lemma}
\begin{proof}
For a permutation $p=(p_1,\ldots,p_r)$ of $1,\ldots,r$, put $H=r+1$ and
\[
 W_p=0\,1^H\prod_{i=1}^{r}(0^i1^{p_i}0^i1^H)\,0^r.
\]
The following nonempty palindromic factors occur in $W_p$:
\[
\begin{array}{ll}
 0^t,\quad 0^t1^H0^t & 1\le t\le r,\\
 0^t1^{p_i}0^t & 1\le i\le r,\quad 1\le t\le i,\\
 0^t1^H0^i1^{p_i}0^i1^H0^t
   & 1\le i\le r,\quad 1\le t\le\max\{1,i-1\},\\
 1^t & 1\le t\le H,\\
 1^t0^i1^t & 1\le i\le r,\quad 1\le t\le p_i,\\
 1^t0^i1^{p_i}0^i1^t & 1\le i\le r,\quad 1\le t\le H.
\end{array}
\]
In the third row, the zero-run before the left $1^H$ has length $1$
when $i=1$ and $i-1$ otherwise; the zero-run on the right is at least
as long. This gives the stated range of $t$.
The run patterns, $p_i<H$, and the distinct values $p_i$ distinguish
all the listed words. Their total number is $L_r=|W_p|$, so $W_p$ is rich.
Its longest zero-run has length $r$; appending one zero adds the new
palindrome $0^{r+1}$. Thus $W_p0$ is rich.
The permutation is recoverable from the successive short one-runs,
and $|W_p|_0=(r+1)^2$. Taking $z$ to be the bitwise complement of $W_p$
gives the assertion.
\end{proof}

\begin{proposition}\label{prop:product}
For $q\ge1$ and positive integers $r_1,\ldots,r_q$, define
\[
 N=1+\sum_{j=1}^{q}\bigl((2j-1)L_{r_j}+(r_j+1)^2\bigr).
\]
Then $R_2(N)\ge\prod_{j=1}^{q}r_j!$.
\end{proposition}
\begin{proof}
Choose the words in Lemma~\ref{lem:permutation} independently and apply
Lemma~\ref{lem:concat}. A zero of block $j$ contributes $2j-1$ letters
to the gap encoding, and a one contributes $2j$, giving the stated
common length. Parsing at zeros recovers the integer word; the value
$\lfloor a/2\rfloor$ identifies the block containing each integer $a$.
Thus every permutation is recoverable, and all outputs are distinct.
\end{proof}

\begin{proof}[Proof of Theorem~\ref{thm:main}]
Fix $k\ge0$. For integers $s\ge2$, take $q=2^{k+1}-1$ and
\[
 r_j=2s\,2^{k-t}\qquad(2^t\le j<2^{t+1},\quad 0\le t\le k),
\]
and denote the length in Proposition~\ref{prop:product} by $N_{s,k}$.
Since $(2u)!\ge s^u$ for $u\ge s$, that proposition gives
\begin{equation}\label{eq:count}
 \log R_2(N_{s,k})\ge (k+1)s2^k\log s.
\end{equation}
Also $L_r\le9r^2$ and $(r+1)^2\le4r^2$ for $r\ge1$.
The $2^t$ blocks at level $t$ have total length contribution at most
\[
 2^t(9\cdot2^{t+2}+4)(2s\,2^{k-t})^2
 \le160s^2 4^k.
\]
The first block also gives $N_{s,k}\ge s$. Thus
\begin{equation}\label{eq:length}
 s\le N_{s,k}\le200(k+1)s^2 4^k+1.
\end{equation}
Letting $s\to\infty$ at fixed $k$ in
\eqref{eq:count}--\eqref{eq:length} gives
\[
 \liminf_{s\to\infty}
 \frac{\log R_2(N_{s,k})}{\sqrt{N_{s,k}}\log N_{s,k}}
 \ge \frac{\sqrt{k+1}}{2\sqrt{200}}.
\]
Since $N_{s,k}\to\infty$ and $k$ is arbitrary, the theorem follows.
\end{proof}

\paragraph{Formal verification.}
Theorem~\ref{thm:main} and the negation of \eqref{eq:conjecture} have
also been formalized in Lean~4 using Mathlib. The formal proof uses the
product construction and explicit integer bounds for the scaled profiles above.
The formalization and paper source are available at\\
\url{https://github.com/glasss13/rich-word-counting}.

\paragraph{Acknowledgment.}
GPT-6 Astra (OpenAI) was used for the main proof construction and drafting
of this paper.

\begin{thebibliography}{9}
\bibitem{GJW}
A.~Glen, J.~Justin, S.~Widmer, and L.~Q.~Zamboni,
Palindromic richness,
\emph{European Journal of Combinatorics} \textbf{30}(2) (2009), 510--531.
\href{https://doi.org/10.1016/j.ejc.2008.04.006}{doi:10.1016/j.ejc.2008.04.006}.

\bibitem{GSSpre}
C.~Guo, J.~Shallit, and A.~M.~Shur,
On the combinatorics of palindromes and antipalindromes,
arXiv preprint \href{https://arxiv.org/abs/1503.09112v1}{arXiv:1503.09112v1} (2015).

\bibitem{GSS}
C.~Guo, J.~Shallit, and A.~M.~Shur,
Palindromic rich words and run-length encodings,
\emph{Information Processing Letters} \textbf{116}(12) (2016), 735--738.
\href{https://doi.org/10.1016/j.ipl.2016.07.001}{doi:10.1016/j.ipl.2016.07.001}.

\bibitem{Ruk}
J.~Rukavicka, Subexponential upper bound on the number of rich words,
arXiv preprint \href{https://arxiv.org/abs/2511.11069v3}{arXiv:2511.11069v3} (2026).
\end{thebibliography}
\end{document}
